Visual moral attribution requires two judgements: whether an image is morally
significant and which moral concern best explains it. Because images often
leave motives and social context unstated, fluent responses can conceal
different moral readings. This dissertation examines those differences rather
than reducing them to one alignment score.

GPT-4o, GPT-5.6, and Qwen3.6 Flash each classified 127 images using ten
categories adapted from Moral Foundations Theory and a \textit{None} option.
Agreement was compared at the exact-label, foundation, and
moral/\textit{None} levels. Further analyses used three repeat runs on 32
images, 280 multi-label judgements from 56 participants on 16 images, and 11
edited variants; alternative prompts provided a supplementary robustness
check.

Agreement depended strongly on resolution. Pairwise agreement rose from
67.7--72.4\% on exact labels to 85.8--88.2\% on the
moral/\textit{None} distinction. Even among the 60 images all three models
considered moral, only 33 received the same exact label. The models therefore
agreed more readily that a scene was morally significant than on what it meant.
Within-block repeatability was high (85.4--95.8\%), but primary-to-repeat
continuity ranged from 62.5\% to 100\%, showing that short-run consistency did
not guarantee continuity across evaluation windows.

Participants also did not provide one answer key: on eight of the 16 survey
images, no category reached 50\% endorsement. GPT-5.6 showed the highest
descriptive participant correspondence, although its lead came entirely from
the six model-disagreement images. Reported confidence was only weakly
associated with participant endorsement. In the edit cases, five of eight label
changes came from two variants that changed the scene's apparent meaning,
whereas obscuring a suspected cue often left the label intact. Alternative
prompts did not consistently improve participant correspondence.

Overall, the models appeared to construct moral readings from visible evidence
and contextual assumptions rather than retrieve a category contained in an
image. Moral relevance, framing, participant correspondence, repeatability,
continuity, and cue sensitivity are distinct properties and should be evaluated
separately.

The findings are methodologically relevant to SDG~16 Targets~16.6--16.7 and,
indirectly, to SDG~10; the study does not measure progress towards either goal.

\keywords{large vision--language models, visual moral attribution, moral
foundations, participant correspondence, model agreement, repeatability,
visual cues}
